\maketitle
\makesignature

\ifproject
\begin{abstractTH}
โครงงานนี้มุ่งศึกษาประสิทธิภาพของแบบจำลองการแบ่งส่วนภาพเชิงความหมาย (semantic segmentation) สำหรับการระบุอวัยวะในช่องปากภายใต้บริบทที่มีรอยโรคร่วม ซึ่งเป็นปัญหาที่มีความท้าทายสูง เนื่องจากรอยโรคสามารถก่อให้เกิดการบิดเบือนลักษณะเชิงกายวิภาค ส่งผลให้แบบจำลองมีความคลาดเคลื่อนในการทำนาย โดยเฉพาะในกรณีข้อมูลที่ไม่ปรากฏในชุดฝึก

การศึกษานี้ดำเนินการประเมินแบบจำลองการเรียนรู้เชิงลึกหลายสถาปัตยกรรม ได้แก่ U-Net, LinkNet, DeepLabV3+ และ UPerNet ร่วมกับ encoder ที่หลากหลาย ภายใต้กรอบการทดลอง 3 รูปแบบ ได้แก่ แบบจำลองพื้นฐาน (baseline), แบบจำลองร่วมกับกระบวนการเตรียมข้อมูล (preprocessing) และแบบจำลองที่ผสาน preprocessing และการปรับปรุงผลลัพธ์ (post-processing)
ในส่วนของ preprocessing ได้มีการออกแบบวิธีการเพื่อปรับปรุงคุณภาพข้อมูลและลดความแปรปรวนของภาพ ขณะที่ post-processing ใช้แนวทางเชิงกฎ (rule-based) และการวิเคราะห์การเชื่อมต่อของพิกเซล (connectivity analysis) เพื่อแก้ไขข้อผิดพลาดเชิงโครงสร้างและเพิ่มความสอดคล้องทางกายวิภาคของผลลัพธ์

ผลการทดลองเชิงปริมาณและเชิงคุณภาพแสดงให้เห็นว่า การผสาน preprocessing และ post-processing สามารถยกระดับประสิทธิภาพของแบบจำลองได้อย่างมีนัยสำคัญ โดยเฉพาะในกรณีภาพที่มีรอยโรค ซึ่งช่วยลดข้อผิดพลาดเชิงพื้นที่และปรับปรุงความสมเหตุสมผลของ segmentation mask
งานวิจัยนี้ชี้ให้เห็นถึงข้อจำกัดของแบบจำลอง segmentation ภายใต้ข้อมูลที่มีความซับซ้อนสูง และเสนอแนวทางในการปรับปรุง pipeline เพื่อเพิ่มความทนทาน (robustness) ของระบบในบริบทการใช้งานจริงทางการแพทย์
\end{abstractTH}

\begin{abstract}
This study investigates the performance of semantic segmentation models for oral tissue segmentation under lesion-aware conditions, which pose significant challenges due to anatomical distortions caused by pathological regions. Such variations often degrade model generalization, particularly for unseen data distributions.

Multiple deep learning architectures, including U-Net, LinkNet, DeepLabV3+, and UPerNet, are systematically evaluated in conjunction with various encoder backbones. The experimental framework consists of three configurations: (1) baseline models, (2) models enhanced with preprocessing, and (3) models integrating both preprocessing and post-processing.
The preprocessing pipeline is designed to improve data quality and reduce inter-sample variability, while the post-processing stage incorporates rule-based refinement and pixel connectivity analysis to correct structural inconsistencies and enforce anatomical plausibility in the segmentation outputs.

Both quantitative and qualitative results demonstrate that the integration of preprocessing and post-processing leads to significant performance improvements, particularly in lesion-containing images. The proposed approach effectively reduces spatial errors and enhances the structural coherence of predicted segmentation masks.
This work highlights the limitations of segmentation models in handling complex medical imaging scenarios and proposes a practical pipeline to improve model robustness and reliability for real-world clinical applications.
\end{abstract}

\iffalse
\begin{dedication}
This document is dedicated to all Chiang Mai University students.

Dedication page is optional.
\end{dedication}
\fi % \iffalse

\begin{acknowledgments}
โครงงานเรื่อง "การแบ่งส่วนเนื้อเยื่อในช่องปากโดยตระหนักถึงการมีของรอยโรค" จะไม่สําเร็จลุล่วงได้ หากไม่ได้รับความกรุณาและความช่วยเหลือจากหน่วยงานต่าง ๆ และท่านผู้มีพระคุณหลายท่าน

ข้าพเจ้าขอกราบขอบพระคุณ รองศาสตราจารย์ ดร.ปฏิเวธ วุฒิสารวัฒนา อาจารย์ที่ปรึกษาโครงงาน จาก
คณะวิศวกรรมศาสตร์ มหาวิทยาลัยเชียงใหม่ ที่ได้เสียสละเวลาให้คําแนะนํา ให้ข้อเสนอแนะ รวมถึงแนวทาง
ในการดําเนินงานโครงงานตลอดระยะเวลาการศึกษา อีกทั้งยังช่วยตรวจสอบ แก้ไขข้อบกพร่องต่าง ๆ ด้วย
ความเอาใจใส่ จนทําให้โครงงานนี้มีความสมบูรณ์ยิ่งขึ้น

ข้าพเจ้าขอขอบคุณ คณะวิศวกรรมศาสตร์ มหาวิทยาลัยเชียงใหม่ ที่ได้เอื้ออํานวยสถานที่ อุปกรณ์ และ
ความสนับสนุนทางด้านต่าง ๆ ที่จําเป็นต่อการทําโครงงาน ตลอดจนการอํานวยความสะดวกในทุกขั้นตอน
ของการดําเนินงาน

ขอขอบพระคุณ ศูนย์ทันตสาธารณสุขระหว่างประเทศ กรมอนามัย กระทรวงสาธารณสุข ที่ให้การสนับสนุนข้อมูลและความรู้ทางด้านทันตสาธารณสุข อันมีประโยชน์อย่างยิ่งต่อโครงงานนี้

ขอขอบคุณ คณะทันตแพทยศาสตร์ มหาวิทยาลัยเชียงใหม่ และ คณะทันตแพทยศาสตร์ มหาวิทยาลัย
ธรรมศาสตร์ ที่ให้ความอนุเคราะห์ด้านข้อมูล ด้านวิชาการ และร่วมสนับสนุนการดําเนินงานของโครงงานใน
ครั้งนี้

ขอขอบคุณผู้ปกครอง เพื่อน และผู้มีส่วนเกี่ยวข้องทุกท่าน ที่ให้กําลังใจ คําแนะนํา และความช่วยเหลือ
ในด้านต่าง ๆ ตลอดระยะเวลาการทําโครงงาน ซึ่งเป็นแรงผลักดันสําคัญที่ทําให้โครงงานฉบับนี้สําเร็จลุล่วงได้
อย่างสมบูรณ์

สุดท้ายนี้หากโครงงานนี้มีข้อผิดพลาดประการใด ข้าพเจ้าขอกราบขออภัยมา ณ ที่นี้ และพร้อมน้อมรับ
ข้อเสนอแนะเพื่อนําไปปรับปรุงให้ดียิ่งขึ้นต่อไป

\acksign{2026}{3}{27}
\end{acknowledgments}%
\fi % \ifproject

\contentspage

\ifproject
\figurelistpage

\tablelistpage
\fi % \ifproject

% \abbrlist % this page is optional

% \symlist % this page is optional

% \preface % this section is optional
